\section{Follow-up 1: the transport law and actual bounded detectors}
\label{sec:f1-transport}

The identifiers in this section belong to DEG\_L Q-BOUNDS and Q-DIV.
We retain the definition
\[
 \delta_W(A,\theta)=\min\{\deg_\theta Z:
       Z\subset A\text{ integral of dimension }n,
       \ e_W[Z]\ne0\},\qquad \dim A=2n.
\]
The minimum of the empty set is $+\infty$.
The supplied B19, B21, B22a and T6.24--T6.25 keep their stated
hypotheses. In particular B19 and B21 use a fixed simple CM reference
with rank-one homology over a CM field; ATTACK Lemma H uses the
commutative algebra $K_0\times K_0$ on $E_0^2\times E_0^2$.
The following actual core morphism gives the comparison needed there.

\subsection{HD-F1-T01: core factorization for the product CM algebra}

Let $(A,\theta)$, with $A=V_{\mathbb R}/\Lambda$, have the fixed
integral maximal $\mathcal O_K$ Weil action of signature $(n,n)$,
with polarization Rosati involution restricting to conjugation on $K$.
Set $q_A(u,v)=(-1)^n\int_Au\wedge v$ on the middle Weil plane.
Let $C$ be a
commutative product of CM fields acting rationally and holomorphically
on $A$, containing that $K$ action. Define
\[
 \Lambda^-_C=\{x\in\Lambda:\mathcal O_Cx\subset\Lambda\},\quad
 J_C=[\Lambda:\Lambda^-_C],\quad
 A^-_C=V_{\mathbb R}/\Lambda^-_C,
 \quad i:A^-_C\longrightarrow A.
\]
The integral $K$ action makes both lattices projective
$\mathcal O_K$-modules. Put
\[
 \mathfrak j_C=\det_{\mathcal O_K}(\Lambda^-_C)
                  \det_{\mathcal O_K}(\Lambda)^{-1},
 \qquad N\mathfrak j_C=J_C.
\]
Every $C$-linear homomorphism $h:B\to A$ from an abelian variety
with integral $\mathcal O_C$ action factors as an actual morphism
\[
 B\xrightarrow{h^-}A^-_C\xrightarrow{i}A.
 \tag{F1T1.1}
\]
Indeed its integral homology image is $\mathcal O_C$-stable and
therefore contained in $\Lambda^-_C$. The universal-cover linear map
then descends to $h^-$; the resulting homomorphism of complex tori is
algebraic because both tori are abelian varieties.

Write $\mathscr W_A=W_A\cap H^{2n}(A,\mathbb Z)$ and choose the
fixed integral projector $P_A=R e_W$. Applying the determinant-lattice
formula already proved in DEG\_L T6.24 to the isogeny $i$ gives
\[
 i_*\mathscr W_{A^-_C}
   =J_C\mathfrak j_C^{-1}\mathscr W_A
   =\overline{\mathfrak j_C}\mathscr W_A.
 \tag{F1T1.2}
\]
No rank-one hypothesis on $C$-homology occurs in this calculation:
the determinant is taken over the specified imaginary quadratic
field $K$ on a rank-$2n$ lattice. On the common rational cohomology
$i^*$ is the identity and $i_*$ is multiplication by $J_C$.
For an integral signed cycle $\Gamma$ whose class is in the integral
span of proper transports $h_*Z$ as in (F1T1.1), one consequently has
\[
 P_A[\Gamma]\in\overline{\mathfrak j_C}\mathscr W_A.
 \tag{F1T1.3}
\]
The projection is applied on $A^-_C$ after $(h^-)_*$; thus no
same-degree projector is commuted through a map from a
different-dimensional source.

Put $\Delta_K=|\operatorname{disc}K|$ and
$V_\theta=\int_A\theta^{2n}/(2n)!$.
For a nonzero projection, T6.24's exact norm-index identity gives
\[
 q_A(P_A[\Gamma],P_A[\Gamma])
     \in 2\Delta_K^{n-1}J_C\mathbb Z_{>0}.
 \tag{F1T1.4}
\]
Using the elementary mass estimate of ATTACK R4, already carried
in HD-D01, yields the fully sufficient bound
\[
 D_{\rm abs}(\Gamma)\ge
 \frac{\sqrt2\,n!\sqrt{V_\theta}\Delta_K^{(n-1)/2}}{R}
           \sqrt{J_C}.
 \tag{F1T1.5}
\]
This displayed constant deliberately uses only that mass estimate;
DEG\_L T6.25 has its stronger calibration constant within its own
stated scope. Equation (F1T1.5) follows by bounding the norm of
$P_A[\Gamma]$ above by
$R D_{\rm abs}(\Gamma)/(n!\sqrt{V_\theta})$ and below by
the square root in (F1T1.4).
For reduced images of generic degrees bounded by $S$, multiplication
by $\operatorname{lcm}(1,\ldots,S)$ puts their integral span into
(F1T1.3), exactly as in T6.25(B). This specifies the actual morphism
and the denominator responsible for the comparison.

\subsection{HD-F1-T02: an eight-degree detector on a specified Lemma H sequence}

Let $K_0=\mathbb Q(\omega)$, $\omega^2+\omega+1=0$, and
$E_0=\mathbb C/\mathcal O_{K_0}$ with its principal polarization.
Start from $A_0=E_0^4=E_0^2\times E_0^2$, with product principal
Riemann form $\psi$ and balanced Weil action
\[
 \iota(a)=\operatorname{diag}(a,a,\bar a,\bar a).
 \tag{F1T2.1}
\]
Fix a rational prime $p=\pi\bar\pi$ split in $K_0$, where
$\pi\in\mathcal O_{K_0}$ and $N\pi=p$; for example
$p=7$, $\pi=3+\omega$.
We now specify the product bases left open in ATTACK Lemma H.
At the place of the balanced action corresponding to $\pi$ take
$e_i$ in the first CM block and $f_i$ in the second, for $i=1,2$.
Their $\psi$-dual basis is $e_i^\vee,f_i^\vee$ in the conjugate
balanced eigenspace. Choose the bases by using identical elliptic
factor bases in both $i$-blocks: $f_i$ is the copy in the second
CM block of the elliptic conjugate-eigenvector $e_i^\vee$.
Then the actual coordinate swap between the two elliptic factors of
one paired block satisfies
\[
 \tau(e_i)=-f_i^\vee,\quad \tau(f_i)=e_i^\vee,
 \quad \tau(e_i^\vee)=f_i,\quad \tau(f_i^\vee)=-e_i.
 \tag{F1T2.2}
\]
These signs express the alternating pairing, not a change of the
global swap, whose square is one.

At $p$ put
\[
 \begin{split}
 M_k&=\bigoplus_{i=1}^2
       \bigl(\mathbb Z_p(e_i+p^{-k}f_i)\oplus\mathbb Z_p f_i\bigr),\\
 M_k^\vee&=\bigoplus_{i=1}^2
       \bigl(\mathbb Z_p e_i^\vee\oplus
                  \mathbb Z_p(f_i^\vee-p^{-k}e_i^\vee)\bigr).
 \end{split}                                                    \tag{F1T2.3}
\]
Let $\Lambda_k$ have localization $M_k\oplus M_k^\vee$ at $p$
and the original product lattice at all other primes.
This is the matched-product realization of the displayed Lemma H
construction. Its Riemann form is integral and unimodular, and the
unchanged complex structure and $\psi$ satisfy the Riemann bilinear
relations. Thus it defines a principally polarized
$(A_k,\theta_k)$ in ATTACK's principal component.

For $C=K_0\times K_0$ acting on the two original CM blocks,
the largest $\mathcal O_C$-stable sublattice is
\[
 \Lambda_k^-=(\Lambda_k\cap V_1)\oplus(\Lambda_k\cap V_2)
            =\pi^k\Lambda_0.
 \tag{F1T2.4}
\]
Here $\pi$ acts by the original central CM scalar in every
elliptic factor, rather than by the balanced action $\iota(\pi)$.
To verify (F1T2.4), the four local intersections are
\[
 \begin{array}{c|cc}
 &V_1&V_2\\ \hline
 M_k&p^k\mathbb Z_p^2e&\mathbb Z_p^2f\\
 M_k^\vee&\mathbb Z_p^2e^\vee&p^k\mathbb Z_p^2f^\vee.
 \end{array}                                                    \tag{F1T2.5}
\]
The original CM $\pi$-eigenspace is spanned by $e_i,f_i^\vee$;
its conjugate is spanned by $e_i^\vee,f_i$.
These are exactly the valuations of $\pi^k\Lambda_0$.
Away from $p$, $\pi$ is a unit. Equality follows from equality
at every localization. Each of the two balanced summands has core
index $p^{2k}$, so
\[
 \boxed{J_W(A_k)=p^{4k}.}                                    \tag{F1T2.6}
\]
In particular the displayed lower bound $p^{2k}$ in Lemma H
omitted the equal contribution from the dual summand.

Equations (F1T2.3) split orthogonally into the identical $i=1,2$
blocks, also away from $p$. Therefore
\[
 (A_k,\theta_k)=(B_k\times B_k,
                     \operatorname{pr}_1^*\theta_{B_k}
                     +\operatorname{pr}_2^*\theta_{B_k}),
 \tag{F1T2.7}
\]
with $B_k$ a principally polarized abelian surface.
The swap in (F1T2.2) preserves its lattice, since
\[
 \tau(e_i+p^{-k}f_i)=-(f_i^\vee-p^{-k}e_i^\vee),\qquad
 \tau(f_i)=e_i^\vee.
\]
It preserves the product Riemann form and the original complex
structure and hence descends to an algebraic polarized involution
$\tau_k$ of $B_k$. It satisfies
$\tau_k\iota(a)=\iota(\bar a)\tau_k$.

The surface
\[
 Y_k=\operatorname{Graph}(\tau_k)\subset B_k\times B_k=A_k
 \tag{F1T2.8}
\]
is smooth and integral. Its graph embedding $j_k$ gives
\[
 j_k^*\theta_k=\theta_{B_k}+\tau_k^*\theta_{B_k}
             =2\theta_{B_k},\qquad
 \deg_{\theta_k}Y_k=4\int_{B_k}\theta_{B_k}^2=8.
 \tag{F1T2.9}
\]
To check detection, write $H^1(B_k,\mathbb C)=U_\sigma\oplus
U_{\bar\sigma}$ for the two $K$-eigenspaces, each of dimension two.
On the $\sigma$ eigenspace of $H^1(A_k)$, pullback along $j_k$ is
\[
 U_\sigma\oplus U_\sigma\longrightarrow
       U_\sigma\oplus U_{\bar\sigma},\qquad
       (u,v)\longmapsto u+\tau_k^*v.
 \tag{F1T2.10}
\]
It is an isomorphism. Its fourth exterior power is nonzero and
integrates nontrivially on $B_k$. Thus $e_W[Y_k]\ne0$.
We have proved
\[
 \boxed{1\le\delta_W(s_k)\le8,\qquad J_W(s_k)=p^{4k}.}
 \tag{F1T2.11}
\]
No equality for the minimum $\delta_W$ is asserted.
For an arbitrary choice of adapted local bases in Lemma H the
core-index computation still gives $p^{4k}$, while (F1T2.8) is
asserted for the matched product bases just specified. The supplied
global principal-component bound applies to all choices.

The same construction gives the exact transport multiplicity.
Multiplication by the original CM scalar $\pi^k$ induces
\[
 h_k:A_0\longrightarrow A_k,\qquad
 \deg h_k=p^{4k},\qquad h_k^*\theta_k=p^k\theta_0.
 \tag{F1T2.12}
\]
By (F1T2.4) it is the core isogeny after the isomorphism
$A_0\simeq A_k^-$. Write $B_0=E_0^2$ in the paired-block
decomposition and $\Gamma_0=\operatorname{Graph}(\tau_0)\subset
B_0^2$. This is a single fixed degree-eight source surface.
The map $h_k$ is $h_{B,k}\times h_{B,k}$, commutes with the
coordinate swap, and $\deg h_{B,k}=p^{2k}$.
Consequently, as integral cycles,
\[
 \boxed{(h_k)_*[\Gamma_0]=p^{2k}[Y_k],\quad
   \deg_{\theta_k}(h_k)_*\Gamma_0=8p^{2k}
                                =8\sqrt{J_W(s_k)},\quad
   \deg_{\theta_k}(h_k\Gamma_0)_{\rm red}=8.}
 \tag{F1T2.13}
\]
This gives the refined statement and a counterexample to the
unqualified assertion in ATTACK Theorem E and the intake that
bounded target degrees are never obtained by transport.
The explicit reduced image has bounded degree even though its
proper push has square-root degree growth. The fixed source here
is $E_0^4$, and the translation stabilizer of $\Gamma_0$ is the
positive-dimensional abelian subvariety $\Gamma_0$ itself. Thus the
simple-source hypotheses giving the finite stabilizer in B21 and
B22a are not satisfied. All parts of their actual statements remain
available.

For a nonzero rational Weil class $t$ with $\delta t$ integral and
$T=\int t^2$, the already proved HD-D03 applied to (F1T2.9) gives
\[
 \boxed{c_t(s_k)\le285120\,\delta\lceil\sqrt T\rceil.}
 \tag{F1T2.14}
\]
This uses $4455\cdot8^2=285120$; it does not repeat the integral
projector or extraction proof.

\subsection{HD-F1-T03: the exact minima for unrestricted reduced images}

For every surjective $K$-linear isogeny $h:B\to A$, define
\[
 \delta_{\rm red,h}(A)=\min\{\deg_\theta h(Z)_{\rm red}:
 Z\subset B\text{ integral of dimension }n,
 e_W[h(Z)_{\rm red}]\ne0\}.
\]
There is the exact equality, with the value $+\infty$ allowed,
\[
 \boxed{\delta_{\rm red,h}(A)=\delta_W(A).}                \tag{F1T3.1}
\]
Indeed every admissible image is an admissible target subvariety.
Conversely, for every integral target subvariety $Y$, every
irreducible component $Z$ of the finite etale pullback $h^{-1}Y$
dominates $Y$ and has $h(Z)_{\rm red}=Y$.
The components dominate because a finite flat morphism has no
irreducible component over a proper closed subset of an integral
base. At their generic points, ordinary cycle pushforward is
\[
 h_*[Z]=\deg(Z/Y)[Y].                                      \tag{F1T3.2}
\]
The Weil projection of $Z$ is nonzero: pushforward on rational
cohomology is invertible and commutes with the same-dimensional
$K$-Weil projector. This proves the equality by the two actual
operations, taking reduced image and taking a component of inverse
image, rather than by changing a degree normalization.

Apply (F1T3.1) to the core $i:A_C^-\to A$. It proves that
unrestricted reduced-image realizations from maximal-order sources
give every intrinsic detector. Applying (F1T1.5) to an individual
component gives the necessary inequality
\[
 \deg(Z/Y)\,\deg_\theta Y\ge
  \frac{\sqrt2\,n!\sqrt{V_\theta}\Delta_K^{(n-1)/2}}{R}
       \sqrt{J_C}.                                        \tag{F1T3.3}
\]
Thus bounded target degree along $J_C\to\infty$ forces the generic
multiplicity of these lifts to be unbounded. Equation (F1T2.13)
supplies an explicit realization with multiplicity exactly
$\sqrt{J_W}$ and a fixed source cycle.

Define $D_{\rm prop,C}$ by minimizing total absolute degree over
nonzero Weil projections in the proper-transport subgroup of
(F1T1.3). Every such signed representative has a detecting
irreducible component, so $\delta_W\le D_{\rm prop,C}$.
On the specific sequence above,
\[
 \frac{\sqrt6}{27}\sqrt{J_W(s_k)}
    \le D_{\rm prop,C}(A_k)\le8\sqrt{J_W(s_k)},\qquad
 \delta_{\rm red,h_k}(A_k)=\delta_W(A_k)\le8.
 \tag{F1T3.4}
\]
The lower constant uses the carried Eisenstein projector $R=54$.
These are inequalities between explicitly related sets of cycle
representatives and the exact pushforward relation (F1T3.2).

\subsection{HD-F1-T04: actual nonmaximal-source realizations}

For every fixed integer $m\ge2$ and every integral detector
$Y\subset A$ of degree $D$, take a component
$Z\subset[m]^{-1}Y$. This is an explicit finite-etale construction
on the same polarized source $(A,\theta)$, and
\[
 [m](Z)_{\rm red}=Y,\qquad
 1\le\deg(Z/Y)\le m^{4n},\qquad
 \deg_\theta Z\le m^{2n}D.                                \tag{F1T4.1}
\]
For the degree bound, the full pullback cycle has cohomology
$[m]^*[Y]=m^{2n}[Y]$, hence degree $m^{2n}D$, and each component
is effective. The underlying source has exactly the target
maximal-order core index. In particular, on every family already
proved in this paper to have $\delta_W\le D_{\mathcal S}$, this
construction supplies bounded-degree nonmaximal sources at its
nonmaximal CM points, with the explicit bound
$m^{2n}D_{\mathcal S}$, and bounded generic degrees.
For the sequence of HD-F1-T02 one may take $m=2$, giving source
degree at most $128$, generic degree at most $256$, and source
core index exactly $p^{4k}$.

The corresponding necessary core comparison also has an actual
lattice map. For a same-dimensional $C$-linear isogeny $g:B\to A$
between abelian varieties carrying the same fixed integral maximal
$\mathcal O_K$ Weil action (their $C$ orders may be nonmaximal),
\[
 g_*H_1(B,\mathbb Z)^-\subseteq H_1(A,\mathbb Z)^-,
 \qquad (\deg g)J_C(B)\ge J_C(A).                         \tag{F1T4.2}
\]
This follows by taking the index of the inclusion of the source
core inside the target core. For a component $Z_B$ mapping with
generic degree one to a detecting $Y\subset A$, the integral
minimum in T6.24 and finite-isogeny Poincare scaling give
\[
 \deg_\theta(Y)^2\ge
       \frac{2\Delta_K^{n-1}(n!)^2V_\theta}{R^2}\deg g,
 \quad
 J_C(B)\ge
       \frac{2\Delta_K^{n-1}(n!)^2V_\theta}
            {R^2\deg_\theta(Y)^2}\,J_C(A).                \tag{F1T4.3}
\]
The source index in (F1T4.1) satisfies this necessary target-sized
core requirement. The exact quotient construction and prime-by-prime
version already proved in DEG\_L B25 remain cited background.

\subsection{Scope of the two DEG\_L rows and the Grothendieck formulations}

For every fixed Weil component with an actual bound already proved
in the restored paper (HD-Q, HD-S, HD-G, HD-N and HD-R at their
listed scopes), intrinsic bounded realizations are present.
Equations (F1T3.1)--(F1T3.3) settle the corresponding unrestricted
reduced-image existence item at its CM points. Equation (F1T4.1)
settles the nonmaximal-source existence item with an explicit source
degree bound. These statements use the actual existing detectors.
They do not obtain a detector on a remaining component by pulling
back an object whose existence is unresolved.
The supplied bounded-detector extraction, HD-D01--D04 and
HD-S15--S16, then gives the stated division bounds on precisely
those components; on the specified Lemma H sequence the sharper
number is (F1T2.14). The general DEG\_L Q-BOUNDS and Q-DIV rows
retain their remaining components recorded in the main ledger.

\paragraph{Framework morphisms: Hilbert and Chow schemes (FGA 221).}
The surfaces $Y_k$ are actual smooth graph subschemes, hence
Hilbert points and effective Chow points. On cycle groups the
proper map induced by $h_k$ sends
$[\Gamma_0]$ to $p^{2k}[Y_k]$. Removing that multiplicity is the
explicit reduced-image operation. For any fixed isogeny, the
surjectivity on integral subvarieties used in (F1T3.1) is furnished
by components of the finite-etale inverse image. These operations
are exactly the ones entering the relative proper Chow spaces
used by the already proved bounded-degree propagation statements.

\paragraph{Variational Hodge classes (Grothendieck 1966).}
The nonzero projected class $e_W[Y_k]$ belongs to the Weil local
system of the ambient four-dimensional principal Weil family.
The entire class $[Y_k]$ also has its specific Hodge locus.
The follow-up fixed-surface calculation computes that locus and
its actual obstruction map; it is needed before extending this
particular graph as a surface. Formula (F1T2.14) is an actual
fibrewise algebraic realization on every point of the stated
sequence. The global algebraic realizations on already settled
families are supplied by the corresponding detector theorems.

\paragraph{Biextensions (SGA 7 I, VII--VIII).}
Let $\lambda_L:A\to\widehat A$ be the homomorphism represented
by the polarization biextension. The degree calculations are
the following exact identities of homomorphisms:
\[
 \widehat h_k\lambda_{\theta_k}h_k=p^k\lambda_{\theta_0},
 \qquad
 \widehat j_k\lambda_{\theta_k}j_k
   =\lambda_{\theta_{B_k}}+
        \widehat\tau_k\lambda_{\theta_{B_k}}\tau_k
   =2\lambda_{\theta_{B_k}}.                             \tag{F1T5.1}
\]
Pulling the Poincare biextension back along these morphisms gives
the polarization classes used in (F1T2.9) and (F1T2.12).

\paragraph{Adams operations (SGA 6 V).}
For the graph embedding $j_k:Y_k\hookrightarrow A_k$, both
tangent bundles and the normal bundle are trivial. Grothendieck--
Riemann--Roch gives the exact rational Chern character
\[
 \operatorname{ch}((j_k)_*\mathcal O_{Y_k})=[Y_k].
\]
Its sole component has codimension two. Hence the Adams operation
on the Grothendieck group of perfect complexes satisfies
\[
 \operatorname{ch}\bigl(\psi^a((j_k)_*\mathcal O_{Y_k})\bigr)
       =a^2[Y_k].                                        \tag{F1T5.2}
\]
Likewise finite-etale pushforward on $\Gamma_0$ and ordinary
Riemann--Roch give
$\operatorname{ch}((h_k)_*\mathcal O_{\Gamma_0})
 =p^{2k}[Y_k]$ as classes on $A_k$.
These are the explicit maps and integer factors relating the
cycle calculation to the SGA 6 operations. No unproved
commutation of Adams operations with integral finite pushforward
is used.

\paragraph{Sources and novelty scope.}
The underlying cycle pushforward and finite-flat pull-push are
Stacks Project, tags 02RF and 02RH; the full short proof of
02RH was read. Milne's \emph{Abelian Varieties}, Section 7,
supplies the standard isogeny terminology. The exact quadratic
integral Weil-lattice identity is the supplied DEG\_L T6.24,
used and rechecked here at its determinant and norm-index steps.
The fixed simple-source arithmetic remains DEG\_L B19/B21/B22a.
The explicit lattice sequence and degree-eight graph calculation
above were independently derived in this run. Literature searches
were made in arXiv, journals and Hexagon; van Geemen's
\emph{Weil Classes and Decomposable Abelian Fourfolds},
SIGMA 18 (2022), 097, is related background on decomposable
fourfolds. No claim that graph cycles or isogeny pull-push are
new is made. The recorded search does not establish priority
for the particular degree/index computation.
